\section{局限性}
\label{sec:limitation}

\paragraph{Token 效率。}
Step 3.5 Flash 已达到前沿级智能，但当前在达到与 Gemini 3.0 Pro 相当质量时仍需要更长的生成轨迹。
下一步我们将对思维过程进行剪枝与压缩，以提升效率并保持相同的竞争性表现。


\paragraph{高效的通用与专精统一。}
我们希望将通用能力与深度领域专长统一。为高效实现这一目标，我们正在推进 on-policy 蒸馏的变体，使模型以更高样本效率内化专家行为。


\paragraph{面向开放世界智能体任务的 RL。}

尽管 Step 3.5 Flash 在学术智能体基准上表现有竞争力，但下一阶段的智能体 AI 需要将 RL 应用于专业工作、先进工程与科学研究中的复杂专家级任务。解决这些挑战是部署具备真正自主能力的智能体的前提。

\paragraph{适用范围与约束。}
Step 3.5 Flash 面向编码与工作型任务优化，但在分布迁移时稳定性可能下降。这通常发生在高度专业化领域或长程多轮对话中，模型可能出现重复推理、中英混杂输出，或在时间与身份认知上产生不一致。